\chapter{Novo capítulo}

\noindent{}\begin{quote}
``Conta o Baal Schem Tov\footnote{Rabi Israel ben Eliezer (1698--1760),
  também conhecido como Baal Shem Tov, foi um dos principais rabinos
  poloneses e está entre os precursores do judaísmo \emph{hassídico}.
  Baal Shem Tov significa em português: ``Um com um bom nome''. O judaísmo \emph{hassídico} é um movimento da mística judaica
  que dá ênfase à alegria religiosa e é organizado em grupos chamados de
  dinastias \emph{hassídicas}, governados por rabinos cuja liderança é
  herdada nas famílias.}. \emph{Dizemos: o Deus de Abraão, o Deus de Isaac e
o Deus de Yaacov, e não o Deus de Abraão, Isaac e Yaacov; pois Isaac e
Yaacov não se apoiavam nos estudos e serviços de Abraão, mas eles mesmos
buscavam a unidade do Criador e o seu serviço}''\footnote{\textsc{buber}, op.
  cit., p.\,95.}.
\end{quote}

A filosofia judaica parte do pressuposto de que, dentro de uma sala de
aula ou dentro de uma mesma família, todos os indivíduos --- sejam eles
estudiosos, historiadores, \emph{tzadikim}\footnote{\emph{Tzadik} (no
  plural: \emph{Tzadikim}) é um título concedido a personalidades do
  judaísmo ortodoxo consideradas santos, mestres e/\,ou líderes religiosos
  e espirituais.}, alunos, alunas e mestres --- são postos em debate e
convidados a trazer um novo contexto, novas interpretações e ideias
sobre os textos sagrados\footnote{Conta Martin Buber: ``Um
  discípulo perguntou ao Maguid de Zlotzchov: --- Está escrito no livro
  de Elias: \emph{Todos em Israel têm o dever de dizer: Quando chegará a
  minha obra aos pés da obra de meus patriarcas Abraão, Isaac e
  Yaacov?}. Como se deve entender isso? Como podemos atrever-nos a
  pensar que seríamos capazes de imitar os patriarcas? Explicou o Rabi:
  \emph{Assim como os patriarcas inventaram novas maneiras de servir, cada
  qual um novo serviço conforme sua peculiaridade, um pela mercê, o
  outro pela força, o terceiro pela glória, assim também nós devemos,
  cada um a seu próprio modo e à luz dos ensinamentos e do serviço,
  inventar inovações e não fazer o que já foi feito, mas o que ainda
  está por fazer}.'' \textsc{buber}, \emph{op.\,cit.}, p.\,190.}. Ao colocar estes diversos indivíduos
em um mesmo espaço --- seja ele físico, como uma sala de aula, ou
temporal, que transcende a linha cronológica --- a noção
temporal-espacial judaica se vincula diretamente às argumentações e
constatações dos documentos e das palavras ali escritas.

Nesta lógica, para além de espaços-tempos retilíneos, cronológicos e
progressivos, a noção judaica de espaço-tempo está relacionada
diretamente aos textos --- e aos livros --- e às suas perenidades de
interpretações. São convocadas a uma sala de estudos --- o
\emph{cheder}\footnote{O \emph{cheder} ocupa um papel central no
  folclore judaico, em especial no mundo iídiche, e simboliza a
  escola primária --- local de excelência da educação judaica, em que
  eram ensinados fundamentos do judaísmo e do hebraico bíblico. Os
  estudos eram realizados normalmente em um quarto na casa do
  \emph{melamed} (professor) e as aulas reuniam crianças das mais
  variadas idades.}--- e ali iniciam um debate as mais diversas gerações
de estudiosos e intérpretes dessa tradição. Moisés é lido ao lado de
Rabi Akiva, assim como Baal Schem Tov ao lado de Rabi Iehiel Mihal --- a
hierarquia temporal é suspensa e as noções são passíveis de
interpretação e de adaptação. E finalmente o aluno é posto ao lado do
mestre e vice-versa, até que ele próprio se torne o mestre.

É a partir do preceito de que --- na tradição judaica --- qualquer
indivíduo é convidado a fazer um \emph{chidush} que se constrói o
entendimento de que, na falta de lugares santos --- especificamente
físicos e ao mesmo tempo judaicos ---, existem milhares de tempos
santos. Como já dizia o rabino americano Abraham Hessel, ``o judaísmo constrói catedrais no tempo''. O conceito de ``lugar
santo'' --- tangível e intangível --- foge de arquiteturas estáveis e
intransportáveis.

Os lugares santos judaicos são os símbolos, os signos, os objetos
rituais, litúrgicos e religiosos que sacralizam e santificam um
território-corpo --- em suas mais variadas escalas --- e delimitam uma
silhueta de pertencimento, mesmo que temporária. Por fim, estes lugares
santos são artifícios utilizados em momentos de culto religioso e de
reza, e vinculados a uma experiência e a uma existência de tempos
religiosos, tais quais as rezas diárias e os festejos semanais e anuais.

Em \emph{Derrida, um egípcio --- o problema da pirâmide judia}, o
filósofo alemão Peter Sloterdijk, a partir da narrativa bíblica,
discorre a respeito da influência de uma política egípcia na construção
de uma identidade, de uma religiosidade e de uma espiritualidade
propriamente judaicas. Não à toa são os egípcios escolhidos como
precursores dessa suposta identidade por meio da aglomeração, da
segregação e da exclusão. Pois, para que judeus e judias pudessem ser
considerados um povo --- e isso se dá por meio do recebimento das Tábuas
da Lei ---, este conglomerado de pessoas escravizadas --- filhos dos
patriarcas Abraão, Isaac e Yaacov e das matriarcas Sarah, Rivka, Rachel
e Leah --- foi posto em liberdade por Moisés, vagueou por quarenta anos
no deserto do \emph{Sinai}\footnote{Após a saída do Egito, o recebimento
  das Tábuas da Lei e logo antes da entrada na Terra Prometida, o povo
  judeu enviou para Canaã doze espiões para investigar e colher
  informações a respeito da nova localidade em que se fixariam. Conta a
  história que dez dos doze espiões, ao voltarem, suplicaram ao povo
  hebreu para que não entrasse na terra, inundada de gigantes, feras e
  bestas. O povo, ao receber estas informações, desesperou-se e clamou
  para Deus que preferiam estar escravizados, em retorno ao Egito. Nesse
  momento, Deus entende que para entregar a liberdade a um povo
  escravizado, esses indivíduos --- a geração que foi ao mesmo tempo
  escravizada e posta em liberdade --- deveriam morrer e dar espaço ao
  nascimento de uma nova geração de pessoas já livres. Assim, é exigido
  de Moisés que conduza o povo em retorno e que vagueie com ele por
  quarenta anos no deserto do \emph{Sinai}.} e finalmente criou para si
uma divindade que ultrapassa os monumentos fixos e intransponíveis de
seus perseguidores. O judaísmo inventou para si um Deus portátil que
habita o texto e o livro.

\begin{quote}
O Egito, dessa maneira, se torna o lugar da pirâmide, evidência
incontornável da materialidade do signo, a representação do divino em
irremovíveis e inabaláveis monumentos de pedra. É o contrário do Deus
inventado pelos judeus, um Deus portátil que habita o livro, um Deus não
mais representado por monumentos intransponíveis (pirâmides) e sim por
documentos (o texto sagrado). Esta nova versão de Deus joga inesperada
luz sobre toda uma problemática ligada a meios, transportabilidades,
migrações. Se não é possível transportar deuses que moram em pirâmides
ou templos colossais, é bem mais fácil transportar um Deus que habita o
texto por mais sagrado que seja\footnote{TELLES, \emph{Posto de
  observação: reverberações psicanalíticas sobre cotidiano, arte e
  literatura}, p.\,345.}.
\end{quote}

Nesse sentido, ``meios, transportabilidades, migrações''\footnote{\emph{Ibidem}.}
adquirem uma importância religiosa até então insuspeita. Ou seja, as
diversas tentativas de construção de identidades e de comunidades
judaicas consolidam e incorporam em si as paisagens, situações e
referências daqueles que os humilham e oprimem. Em parte, é a partir de
geografias repressivas que a identidade e as comunidades judaicas
inventam, criam e constroem para si espaços sacralizados --- em suas
múltiplas escalas --- que ultrapassam seus perseguidores e, finalmente,
que se consolidam e florescem.

\begin{center}\adforn{49}\end{center}

\noindent{}A incorporação e a consolidação de traços identitários-comunitários, a
partir de geografias de pertencimento/\,exclusão, podem ser observadas em
meio ao surgimento dos primeiros guetos judaicos na Veneza
renascentista, durante o século \textsc{xvi}. Pois, como nos lembra o historiador
e sociólogo Richard Sennett:

\begin{quote}
Ao acuarem os judeus para o gueto, os venezianos acreditavam estar
isolando o mal que infectara a comunidade cristã. Eles sentiam medo de
tocar os corpos impuros\footnote{Na Veneza renascentista, a comunidade
  cristã acreditava que judeus e meretrizes eram portadores de corpos
  impuros, em que haveria uma prostituição do corpo-carne em nome da
  usura e do dinheiro. Dessa maneira, ambos os grupos sociais eram
  considerados inferiores à moral cristã, sujeitos à exclusão e
  discriminação. Essa relação também é posta em debate ao se pensar a situação
  geográfica de judeus e prostitutas no caso da zona do meretrício no
  bairro do Bom Retiro, na cidade de São Paulo no início do século \textsc{xx}. A
  delimitação de uma determinada zona da cidade em que haveria o
  trânsito de judeus recém-imigrados e prostitutas judias --- as polacas
  --- tensiona a criação de uma narrativa moral para e pela própria
  comunidade judaica. Essa narrativa --- em constante disputa pelas
  comunidades --- exclui a presença das prostitutas judias do bairro e
  da história comunitária, e até da cidade de São Paulo.} que
identificavam como vícios corruptores\footnote{\textsc{sennett}, \emph{Carne e
  Pedra}, p.\,222.}.
\end{quote}

A montagem do gueto se deu no momento em que a cidade de Veneza passava
por um período de recessão econômica. Em consequência, os cristãos
acreditavam que a segregação da diferença --- judeus e demais corpos
indesejáveis, porém necessários --- devolveria à cidade a paz divina e a
dignidade\footnote{As crenças antissemitas que circulavam na época
  sustentavam que os corpos judaicos eram portadores de doenças e de
  vícios --- culpados que eram pela morte de Jesus Cristo e frutos do
  pecado e de impurezas. A comunidade cristã acreditava que a presença
  de judeus --- dos quais muitos praticariam a agiotagem --- teria
  desencadeado uma fúria divina sobre a cidade de Veneza, causando,
  desta forma, sua recessão econômica.}. Dessa maneira, pessoas judias
acabaram sendo restritas a uma geografia repressiva --- o gueto judeu
cercado, onde a vida judaica acontecia. Apesar da construção deste
território plástico\footnote{O conceito de território plástico foi
  desenvolvido pelo pesquisador israelense Eyal Weizman no livro
  \emph{Hollow Land}. Refere-se ao artifício de criação de
  territorialidades artificiais e impositivas perante uma geografia já
  pré-estabelecida. O termo designa as fronteiras fluidas e de
  características plásticas, ou seja, que podem se alargar ou diminuir,
  inflar ou desinchar conforme as necessidades, o contexto e a situação
  em que se encontram.} --- uma territorialidade segregada e excludente
---, a população judaica de Veneza ainda foi capaz de ``construir
novas formas de vida comunitária que extrapolaram as políticas de
isolamento e de conquistar algum nível de
autodeterminação''\footnote{\textsc{sennett}, \emph{op.\,cit.}, p.\,223.}.

Ao final, a geografia de Veneza protegeu os judeus. A construção desta
territorialidade pelos canais fez com que o gueto ficasse isolado,
permitindo que os judeus --- e a vida e a identidade judaicas --- fossem
livres para realizar seus ritos, cultos, tradições e festividades; para
construir para si uma comunidade efervescente por meio da consolidação
de uma identidade compartilhada e de espaços de convívio comunitário
(como os centros comunitários e as sinagogas). Foi a partir de uma
geografia repressiva que os judeus transformaram esta
circunstância-território em uma possibilidade.

Dessa maneira, entende-se que a vivacidade das comunidades judaicas
dentro dos guetos --- seja o gueto de Veneza, o gueto de Varsóvia ou de
diversas outras geografias excludentes ---, durante os momentos de
celebração, de perseguição e de lamento, nos espaços de trânsito, se dá
a partir de constantes movimentos de transformação do espaço maldito ---
o espaço da perseguição --- em um lugar sagrado (\emph{kadosh}),
passível de apropriação, de cultivo das fagulhas pousadas e de
identidades.

Na língua hebraica, a palavra \emph{kadosh} possui dois significados. O
primeiro expressa literalmente ``separar'' ou ``apartado'', o segundo
conota sacralidade, o ``espaço e/\,ou objeto sagrado''. \emph{Kadosh},
logo, simboliza a prática judaica de resguardar momentos e objetos que
são preservados para situações extraordinárias ou de grande valor
simbólico e religioso. A ação de transformar ``\emph{arur}'' (maldito)
em ``\emph{kadosh}'' (sagrado) diz respeito a uma apropriação das
condições e situações em que alguém se vê e, assim, à criação de outros
lugares. Trata-se de decidir por si os próprios passos. \emph{Ansina,
estonses, entiendese ke, endesparte de deshar la luz pasar, es menester
de emitir luz propria}\footnote{Tradução do ladino: ``Isso, pois,
  entende-se que, para além de deixar a luz passar, é necessário emitir
  luz própria''.}.

Por fim, estes movimentos de transitoriedades, transportabilidades e
locomoções --- as constantes apropriações de territórios (geográficos ou
não) --- são observados nos monumentos judaicos e, concomitantemente,
moldam para si os conceitos de identidade e de comunidade --- bem como a
manutenção geracional destes vínculos.

\begin{center}\adforn{49}\end{center}

\noindent{}Contam os \emph{tzadikim} que uma sinagoga nunca deveria ser destruída.
Seu desmonte apenas é possível caso a comunidade que ali frequenta se
locomova --- construindo para si outro espaço comunitário, em outras
localidades. Este preceito ``joga inesperada luz''\footnote{Tomo de
  empréstimo as palavras de Sérgio Telles, \emph{op.\,cit.}, p.\,345.} no
entendimento de que, para além de arquiteturas fixas, os espaços-tempos
judaicos são construídos e consolidados por vínculos de pertencimento,
tradições, línguas, ritos e cultos que resultam em um aglomerado de
indivíduos que compartilham signos e símbolos semelhantes e se agregam,
se unem, constroem e inventam para si uma comunidade.

Ainda nesta chave, não há restrições religiosas para a construção de
espaços sinagogais --- estes locais podem ser construídos e consolidados
em qualquer localidade, contanto que possuam em seu interior as peças
primordiais para o rito religioso. Dentre elas estão: \emph{Aron
HaKodesh}\footnote{\emph{Aron HaKodesh} (em português: ``Arca Sagrada'')
  é uma espécie de armário em que se guarda os textos bíblicos e os
  pergaminhos sagrados, a Torá. Este é um objeto posicionado
  sempre em direção ao Templo de Jerusalém.}, \emph{Bimá}\footnote{\emph{Bimá}
  (em português: ``Altar'') é uma mesa ao centro da sinagoga, usada para
  o suporte de leitura de livros e de textos bíblicos, além de ser o
  local onde o \emph{Chazan} (condutor da oração) se posiciona.} e
\emph{Esh}\footnote{\emph{Esh} (em português: ``Fogo celestial'') é uma
  vela, um candelabro ou uma luz que representa a presença divina na
  casa de orações. Assim como era feito no Templo Sagrado e lembrado durante a
  festividade de \emph{Chanucá}.}. Estes objetos, de cunho religioso,
para além de delimitar uma espacialidade judaica, são móveis,
transitórios e, dessa forma, ocupam, consagram e sacralizam lugares ---
em suas mais variadas escalas ---, animam estas territorialidades,
instalam nelas uma alma.

Por fim, os lugares santos judaicos, aqui investigados, se manifestam na
escala do corpo (\emph{kipá}), da casa (\emph{mezuzá}), da cidade
(\emph{eruv}) e do cosmos (\emph{makom}); cosmogonizam um território
(real, ficcional e/\,ou imaginário) e, para além de lembretes diários,
firmam-se enquanto declarações públicas, individuais, comunitárias de
identidade, de crença e de fé.

\begin{center}\adforn{49}\end{center}

\noindent{}A \emph{kipá}\footnote{A \emph{kipá} é um chapéu, boina ou touca usado
  tradicionalmente por homens e recentemente incorporado por mulheres
  nos movimentos religiosos judaicos reformistas e reconstrutivistas. É
  um objeto que delimita um território-corpo de pertencimento e se firma
  enquanto um lembrete constante da presença de Deus.} e a
\emph{mezuzá}\footnote{A \emph{mezuzá} é um cilindro de metal, de
  madeira ou de acrílico posto nos batentes das portas judaicas e que
  contém um pergaminho com a bênção do \emph{Shemá Israel}. Este ato é
  uma tradição e um lembrete que remete aos tempos de escravização no
  Egito, tendo sido reinterpretado pelas comunidades. Judeus
  sefaraditas costumam posicionar a \emph{mezuzá} de forma
  paralela ao batente, enquanto judeus ashkenazitas a posicionam
  de forma inclinada.} são eixos cósmicos, signos de ligação e de
lembrança da presença divina e de sua conexão com o mundo humano. Estes
símbolos são objetos hierofânicos --- manifestações da realidade sagrada
de Deus por meio de artifícios e de signos do mundo real, e assim
passíveis de compreensão pelas interpretações humanas. Eles se
consolidam enquanto atos misteriosos e distinguem a realidade sagrada de
um mundo profano. Pois ``para o homem religioso o espaço não é
homogêneo\footnote{Em \emph{A poética do espaço}, o filósofo e poeta
  francês Gaston Bachelard aponta que ``nem o tempo, nem o espaço
  em que vivemos são vazios, neutros e homogêneos, mas são carregados de
  qualidade e habitados por fantasmas: os espaços de nossos devaneios e
  paixões podem ser leves, etéreos, transparentes, obscuros ou
  pedregosos, móveis e fixos'', enfim, vivemos em um conjunto de
  relações que, segundo Bachelard, definem nossas experiências,
  agregações individuais e coletivas, e suas possibilidades de
  reinvenção. \textsc{bachelard} \emph{apud} \textsc{rago}, \emph{Inventar outros espaços, criar
  subjetividades libertárias}, p.\,14.}, o espaço apresenta
roturas, quebras; há porções de espaço qualitativamente diferentes de
outros''\footnote{\textsc{eliade}, \emph{O sagrado e o profano}, p.\,25.}.

Neste sentido, a presença divina transcende o mundo espiritual --- em
direção ao mundo material/\,humano --- por meio de objetos humanos
conhecidos e que fazem parte de suas próprias realidades. Por exemplo,
no momento em que a divindade revelou-se a Moisés a partir de uma árvore
em chamas, o objeto ``árvore'' não é adorado como árvore em si, mas sim
enquanto uma hierofania --- um dispositivo que permite a passagem do
sagrado ao mundo humano-profano, já que, ao final, o sagrado se
manifesta quando revela-se.

\begin{quote}
Não te aproximes daqui, disse o Senhor a Moisés; tira as sandálias de
teus pés. Porque o lugar onde te encontras é uma terra Santa\footnote{ÊXODOS,
  3:5.}.
\end{quote}

Assim sendo, ``toda hierofania espacial ou toda consagração de um
espaço equivale a uma cosmogonia''\footnote{\textsc{eliade}, \emph{op.\,cit.}, p.\,59.}.
A \emph{kipá} --- este eixo cósmico de ligação entre o ser humano e o
divino ---, neste sentido, revela-se enquanto consagração de um
determinado território-corpo, um objeto hierofânico portátil, usado por
judeus e por judias como revelação e recordação da presença divina no
mundo --- em que, para além de suas próprias silhuetas, está firmada a
divindade, criadora do mundo e de todas as coisas que ele contém. O eixo
cósmico --- a ligação direta com Deus --- é reforçado a partir da
utilização deste artifício-objeto, pois, ao final, segundo o escritor e
teólogo Mircea Eliade\footnote{\emph{Ibidem}, p.\,35.}, o ``homem
religioso'' deseja estar a todo instante em ``Seu Mundo'', em contato
direto com o divino.

Para além de um objeto hierofânico, a \emph{kipá} se apresenta enquanto
um marcador legal, social, identitário e comunitário de pertencimento
judaico. A utilização deste artifício-objeto se estabelece enquanto um
delimitador de pertencimento e de congregação entre os diversos
indivíduos que transitam em um determinado território onde a comunidade
se estabelece; tais indivíduos, por sua vez, se reconhecem como
portadores de uma identidade compartilhada. E assim, essa peça delimita,
demarca e consagra um corpo/\,\emph{corpus}, possibilitando um senso de agregação
entre os transeuntes de uma determinada territorialidade-cidade onde a
comunidade decide fixar-se.

O uso da \emph{kipá} firma-se enquanto uma possibilidade de
``reconhecer-se judeu ou judia'' em comunidade --- para além do próprio
corpo e da conexão religiosa com o divino. Permanecer em comunidade
permite ao ``homem religioso'' --- este \emph{corpus} de pessoas ---
estar a todo instante em ``Seu Mundo'' --- ou seja, em contato direto
com Deus e com seus iguais. Dessa maneira, o ``homem religioso'', ao
utilizar o artifício dos objetos hierofânicos, cosmogoniza um território
comunitário e cria, por fim, um corpo, um local e um mundo santificado
para si.

A cosmogonização e a consagração de um mundo --- o ``Seu Mundo'' --- são
compreendidas a partir do conceito de eixo cósmico --- uma relação
direta e vertical entre o ser humano e o divino. Segundo Eliade, é a
partir do eixo cósmico --- e em volta dele --- que o território se torna
habitável pela presença divina. Ou seja, é a partir da consagração
destas territorialidades que o ``homem religioso'' faz do espaço profano
--- amorfo e sem limites --- o ``Seu Mundo'', pois ``Nosso Mundo é
a `terra santa' porque é o lugar mais próximo de Deus''\footnote{\emph{Ibidem},
  p.\,40.}.

``A instalação de um território equivale à fundação de um
Mundo''\footnote{\emph{Ibidem}, p.\,46.}. Essa instalação --- ligada ao eixo
cósmico --- ritualiza, sacraliza e demarca um espaço, em suas mais
variadas composições. Parte-se do pressuposto de que o ``homem
religioso'' necessita estar a todo instante em contato com Deus; assim
sendo --- para além de seus corpos ---, estes indivíduos desejam que
suas próprias casas e moradas se tornem o ``Centro do Mundo''. A morada
--- a habitação do ser humano --- situa-se no centro de seu próprio
mundo religioso e reproduz, em escala microcósmica, o Universo criado
por Deus.

Ao utilizar o artifício-objeto hierofânico da \emph{mezuzá}, famílias
judias (e as comunidades) reafirmaram a presença de Deus em suas
próprias moradas. A casa e/\,ou o espaço comunitário (sinagoga), por fim,
se tornam uma \emph{imago mundi}\footnote{\emph{Imago mundi} é um termo
  do latim que representa a ``imagem/\,espelhamento do mundo criado por
  Deus''.} --- ou seja, ao passar por um ritual de consagração e
sacralização, tornam-se um espaço religiosamente habitável, em conexão
com Deus a partir do eixo cósmico. Nos textos talmúdicos, a casa
judaica religiosa é referida como \emph{Bait Neeman} --- uma ``casa
leal'' --- uma habitação santificada pela presença divina, já que ela
própria reitera a criação do mundo e a fidelidade e dedicação ao
monoteísmo e ao Deus judaico.

Essa tradição remete ao momento em que Moisés suplica ao Faraó a
liberdade de seu povo escravizado no Egito. Segundo a narrativa bíblica,
a restrição dessa liberdade impõe que Deus envie ao Egito dez pragas,
sendo a última delas a descida à terra do Anjo da Morte e o assassinato
de primogênitos. Para que os hebreus fossem poupados, Moisés aconselha o
povo para que pinte os batentes das portas com sangue de carneiro. Esse
ritual consagra e delimita um território de crença e de devoção a Deus
--- santifica uma habitação, a torna ``fiel'' à presença divina.

Enquanto uma reinterpretação do sangue de carneiro despejado pelos
hebreus nos umbrais de suas portas, a \emph{mezuzá} é um artifício de
demarcação espacial que eleva um território-casa. É uma caixa de
acrílico, metal ou madeira fixada nos batentes das portas judaicas e que
contém um pergaminho com a bênção do \emph{Shemá Israel}\footnote{A reza
  de \emph{Shemá Israel} é uma das mais importantes bênçãos judaicas;
  afirma a lealdade a Deus e ao monoteísmo.}. Este objeto é portátil,
removível, ajustado e acompanha as comunidades nos momentos de diásporas
e dispersões, voluntárias ou forçadas. A \emph{mezuzá} reitera a
cosmogonia e a devoção judaica de lealdade a Deus, santifica uma
habitação e se estabelece enquanto um simbolismo cósmico, um reflexo da
própria criação divina do mundo. Cria-se, a partir dela, o eixo cósmico.
Já que esse objeto é a chave da casa, fixar-se em um local é instalar
nele, quase que instantaneamente, uma \emph{mezuzá}.

Se toda criação humana repete a fundação do mundo, para que esta criação
dure --- perdure no espaço e no tempo --- ela necessita ser animada,
receber vida, ``uma alma''\footnote{\textsc{eliade}, \emph{op.\,cit.}, p.\,53.}. Segundo
Eliade\footnote{\emph{Ibidem}, p.\,11.}, ``para o homem primitivo tudo é
dotado de alma''. Em uma sinagoga, em um centro religioso ou
comunitário, a comunidade que ali transita é a alma e a chama de vida
dessa localidade (o \emph{minian}). A comunidade --- essa congregação de
indivíduos em prol de objetivos em comum --- funda um território de
pertencimento que anima o lugar --- introduz nele uma alma. Assim sendo,
as comunidades judaicas transformam os espaços profanos em lugares
sagrados por meio da consagração destes territórios. Nestas
territorialidades tudo é dotado de vida, de celebração e de rito, tudo
parte de uma orientação. O ``homem religioso'', por fim, transforma o
maldito em sagrado --- constrói suas identidades, santifica esses
espaços e os compartilha.

Instalar-se em um determinado lugar é reiterar a cosmogonia --- imitando
a obra divina de criação do mundo por Deus. ``Fixar-se em uma
região é assumir responsabilidade por este mundo criado; mantê-lo vivo e
renová-lo''\footnote{\emph{Ibidem}, p.\,54.}. Fixar-se em um determinado local
equivale a um novo começo. E, assim como um calendário cíclico, todo
começo repete o início primordial da criação divina do mundo.

Já que instalar-se em uma determinada territorialidade funda de novo um
mundo, essas tentativas de cosmogonização são formas de dedicar-se a
viver o cosmos puro e santo, tal como era no princípio, quando saiu das
mãos da divindade. Estas propostas são observadas nos objetos
hierofânicos --- \emph{kipá} e \emph{mezuzá} --- e relacionam-se à
experiência de tempos santos judaicos, cíclicos, elipsoidais e
constantes. A vitória do lugar sagrado sobre o espaço profano deve ser
repetida simbolicamente a todo instante, pois a todo instante o mundo
necessita ser criado novamente desde o momento primordial.

Ao fim e ao cabo, a experiência do sagrado funda ontologicamente o
mundo. ``O espaço sagrado tem um valor existencial para o homem
religioso. Nada se pode começar, nada se pode fazer sem uma orientação
prévia''\footnote{\emph{Ibidem}, p.\,26.}. Este ponto fixo --- a criação do
mundo --- permite uma direção. Se, para viver no mundo, o ``homem
religioso'' necessita fundá-lo, nenhum mundo pode nascer do ``caos'', da
homogeneidade e da relatividade do espaço profano.

\begin{center}\adforn{49}\end{center}

\noindent{}Fundar um mundo, orientar-se em meio ao ``caos'' --- ao relativismo do
espaço profano ---, torná-lo sagrado, é sempre possível. Para isso pegue
o primeiro texto que encontrar, \emph{texto-memória},
\emph{texto-qualquer}, um \emph{texto-nenhum}. Ele sempre será um mapa.

Estes monumentos móveis são Deus, que habita o texto e o livro. São
faróis e dispositivos de direção que acompanham as comunidades nos
constantes movimentos das diásporas --- pousam e fixam-se em outras
territorialidades; inventam e fundam outros, mas nunca novos, mundos.